\section{Readout results}
\label{sec:readout}

\begin{table}[t]
\centering\scriptsize
\setlength{\tabcolsep}{3pt}
\begin{tabular}{lccccccccc}
\toprule
& \multicolumn{3}{c}{Qwen2.5-7B} & \multicolumn{3}{c}{Llama-3.1-8B} & \multicolumn{3}{c}{Mistral-7B} \\
\cmidrule(lr){2-4}\cmidrule(lr){5-7}\cmidrule(lr){8-10}
Task & AUROC [95\% CI] & nuis. & gain & AUROC [95\% CI] & nuis. & gain & AUROC [95\% CI] & nuis. & gain \\
\midrule
Sufficiency & 0.948 [0.94,0.95] & 0.87 & 0.11 & 0.952 [0.95,0.96] & 0.86 & 0.12 & 0.949 [0.94,0.95] & 0.87 & 0.11 \\
Sufficiency (matched) & 0.911 [0.90,0.92] & 0.90 & 0.07 & 0.918 [0.91,0.93] & 0.89 & 0.08 & 0.910 [0.90,0.92] & 0.90 & 0.07 \\
Uptake & 0.775 [0.74,0.80] & 0.66 & 0.14 & 0.783 [0.75,0.81] & 0.58 & 0.21 & 0.771 [0.74,0.80] & 0.59 & 0.18 \\
Correction & 0.730 [0.70,0.76] & 0.81 & 0.01 & 0.817 [0.79,0.84] & 0.83 & 0.06 & 0.792 [0.77,0.81] & 0.85 & 0.03 \\
Stability & 0.750 [0.72,0.78] & 0.92 & $-$0.01 & 0.869 [0.85,0.89] & 0.90 & 0.02 & 0.800 [0.77,0.83] & 0.90 & 0.01 \\
State sufficiency & 0.992 [0.99,0.99] & -- & -- & 0.994 [0.99,0.99] & -- & -- & 0.996 [1.00,1.00] & -- & -- \\
\bottomrule
\end{tabular}
\caption{Held-out AUROC of the best-layer probe on paired activation updates (document-disjoint split; 95\% question-level bootstrap intervals), with the AUROC of a classifier on nuisance covariates alone and the gain from adding the probe score to that classifier. The last row is a different task: the state-sufficiency probe classifies single contexts (raw activation $h(C)$, sufficient vs.\ insufficient conditions), not transitions, and its AUROC is not comparable with the rows above (\cref{sec:method,sec:what}).}
\label{tab:probes}
\end{table}

\begin{figure}[t]
\centering
\includegraphics[width=0.84\linewidth]{auroc_by_layer}
\caption{Held-out AUROC by relative depth for the update tasks (document split). Sufficiency peaks around mid-depth in every model; uptake and correction emerge later.}
\label{fig:layers}
\end{figure}

\paragraph{Sufficiency is linearly readable and robust.}
\Cref{tab:probes} shows that the decisive increment is separated from all non-decisive increments at $0.948$--$0.952$ AUROC in the three models, with the best layer at 50--60\% depth (\cref{fig:layers}). Against the matched controls only (distractor, redundant, false, answer-string from the same penultimate state) the probe still reaches $0.91$--$0.92$. The direction survives residualising the nuisance covariates ($0.94$--$0.96$) and adds $+0.11$ AUROC to a nuisance-only classifier. \Cref{tab:splits} shows the result is insensitive to the grouping: question-, document-, answer-entity- and relation-disjoint splits all give $0.94$--$0.96$. The state-sufficiency probe, which classifies whole contexts rather than transitions, reaches $0.99$, but part of that comes from paragraph count, and \cref{sec:what} separates the two. Projections of held-out increments on the sufficiency direction order the increment kinds exactly as hypothesised (\cref{fig:kinds}): decisive $>$ false $>$ distractor $>$ redundant $>$ answer-string control $>$ non-decisive gold, with post-sufficiency increments and the first hop far below. The false-evidence increment scoring second indicates that the direction encodes ``the final slot has been filled'' rather than truth, which is what the model can know before generation.

\begin{figure}[t]
\centering
\includegraphics[width=0.42\linewidth]{projection_by_kind}
\caption{Mean projection of held-out increments on the sufficiency direction by kind (Qwen2.5-7B, document split; standard errors). Red: decisive.}
\label{fig:kinds}
\end{figure}

\paragraph{Uptake is a separate, later signal.}
Whether a decisive increment is actually integrated (answer becomes correct) versus behaviourally ignored is decodable at $0.77$--$0.78$ AUROC (document split; $0.85$ question split), at layers near the end of the network, and this is the task where the activation adds most beyond nuisance covariates ($+0.14$ to $+0.21$). The uptake direction is nearly orthogonal to the sufficiency direction (cosine $0.02$--$0.03$) and the sufficiency direction reaches only $0.54$--$0.60$ on the uptake task, while uptake and wrong-to-correct \emph{correction} share a direction (cosine $0.30$--$0.34$; cross-task AUROC $0.82$; \cref{fig:transfer}, \cref{app:transfer}). Correction and stability, by contrast, are largely explained by the change in the output distribution: nuisance-only classifiers reach $0.81$--$0.85$ and $0.90$--$0.92$, and the probe adds at most $0.06$.

\paragraph{Transfer.}
Applying the 2-hop stateless directions without retraining (\cref{tab:transfer}, \cref{app:transfer}): To the conversational protocol they transfer at $0.79$ (sufficiency), $0.78$ (uptake), $0.92$ (correction) and $0.82$ (stability), supporting the view that the transition is intrinsic to evidence integration rather than an artefact of being asked to reconsider. To \twowiki{} the sufficiency direction transfers at $0.96$ ($0.93$ matched), correction at $0.85$ and uptake at $0.79$; probes trained within \twowiki{} reach $0.99$, $0.89$ and $0.91$. To 3--4-hop questions sufficiency transfers at $0.82$--$0.85$; within that run the leave-one-out completions score as high as chain completions (\cref{app:multihop}), and sufficiency probes reach $0.93$ on every split, but uptake falls to $0.61$ on the document and answer splits (from $0.80$ on the question and relation splits).

\subsection{What the probes read: state, update and output}
\label{sec:what}
Two alternative explanations could account for a high sufficiency AUROC on paired updates. The first is that the update is incidental: the readout may already be present in the absolute state after (or before) the new evidence. The second is that the readout is just the model's planned output. The anchor is the position that emits the first answer token, and the prompt makes INSUFFICIENT the alternative to answering, so a sufficiency probe could be reading ``about to say INSUFFICIENT'' rather than any comparison of evidence with question. We test both on the same increments, labels and splits.

\paragraph{State versus update.}
Refitting every update probe on four representations of the same increments, $h(C_{k-1})$, $h(C_k)$, $\Delta h$ and $[h(C_{k-1});h(C_k)]$ (\cref{tab:rep_ablation}, \cref{app:rep}), shows that most of the sufficiency readout is a property of the state \emph{after} the new evidence. $h(C_k)$ alone reaches $0.930$--$0.940$ against $0.947$--$0.952$ for $\Delta h$ (matched controls: $0.900$--$0.908$ against $0.910$--$0.917$). The state before the evidence reaches $0.733$ on the full task in all three models. This is exactly the AUROC of an oracle that knows only the stage the increment starts from (empty, partial or already sufficient): decisive increments make up $20\%$ of increments that start from a partial state and none of the others, so $h(C_{k-1})$ contributes nothing beyond that stage. On the matched task, where the decisive increment and its four controls share their starting state, it scores exactly $0.500$. Concatenating both states matches $\Delta h$ ($0.949$--$0.950$), so the subtraction adds no information beyond having both endpoints; it removes question-specific content that both states share. For uptake, correction, stability and revision the state after the evidence is as good as or better than the update (uptake $0.769$--$0.798$ against $0.771$--$0.783$; stability $0.813$--$0.882$ against $0.750$--$0.869$). The state-sufficiency probe itself, applied to $h(C_k)$ of the same increments, scores only $0.70$ on the full task: it cannot reject the $27\%$ of negatives that add evidence to an already sufficient context. The task the update probes solve is therefore ``has the context \emph{just} become sufficient'', not ``is it sufficient''. The paired-update design is therefore what makes the comparison controlled: matched controls from the same starting state, and the causal tests of \cref{sec:causal}. It is not what makes the information readable. The model's state represents whether its context has \emph{just} become sufficient, and the update is a clean way to read that.

\paragraph{Evidence versus output.}
The model's output is itself a sufficiency judgement: it answers when it judges the evidence sufficient and says INSUFFICIENT otherwise. That judgement is often wrong on sufficient evidence. On held-out contexts that contain every gold hop, the model says INSUFFICIENT $33\%$ (Qwen), $40\%$ (Llama) and $43\%$ (Mistral) of the time (\cref{tab:behaviour}). As a sufficiency classifier, the verbalised output reaches only $0.68$--$0.76$ AUROC. The state-sufficiency probe, read at the same anchor of the same forward pass, reaches $0.99$. That number is inflated, because every context with three or more paragraphs is sufficient, but with paragraph count held fixed (minimal sufficient contexts against all other two-paragraph contexts, probe retrained on them) it is still $0.976$--$0.986$ against $0.71$--$0.77$ for the verbalised judgement (\cref{tab:verbal_internal}, \cref{app:behaviour}). Restricted to the contexts on which the model says INSUFFICIENT, the probe still separates sufficient from insufficient ones at $0.967$--$0.982$. At $0.95$ training precision it recovers $76$--$81\%$ of the sufficient contexts the model wrongly rejected, with $3$--$4\%$ false alarms on the insufficient contexts it correctly rejected. The sufficiency update probes behave the same way when the output is held fixed. Trained and tested only on increments after which the model \emph{still} says INSUFFICIENT, they reach $0.935$--$0.944$, against $0.65$--$0.69$ for the output alone (\cref{tab:behaviour_control}). The sufficiency readout is therefore not the planned output: the residual stream registers that the evidence has become complete in cases where the model's own answer denies it. Those cases are weaker, though. Wrongly rejected contexts sit $0.80$--$0.87$ of the way from the insufficient to the sufficient class mean, against $1.03$ for answered ones, so the internal state is graded and the verbal decision is a thresholded, noisier read-out of it. This is the gap that uptake measures: a sufficiency signal the output does not act on. It is also why better sufficiency detection does not by itself buy accuracy in adaptive retrieval (\cref{app:adaptive}): stopping on a sufficient context the model will not answer from still yields INSUFFICIENT.

\subsection{Is the internal state needed? Text-only baselines}
\label{sec:text}
\Cref{tab:text} answers the obvious objection. For sufficiency, TF-IDF reaches $0.62$ and a DeBERTa that sees only the question and the added paragraph $0.72$--$0.73$, but a DeBERTa that also sees the previous context reaches $0.89$--$0.92$ at a reduced training budget (one epoch, 256 tokens, 8k increments) and $0.97$ at full budget (two epochs, 512 tokens, all increments), exceeding the probe's $0.95$ on the same split. Sufficiency is therefore largely recoverable from the text: the activation probe is an accurate readout of a property the input already determines. Uptake is the opposite case. On the same increments no text model exceeds $0.63$ on whether the model will integrate the decisive paragraph (Qwen $0.56$ reduced and $0.63$ full budget, Llama $0.52$, Mistral $0.60$ with context), while the activation probes reach $0.77$--$0.78$: whether evidence is taken up depends on the model's internal state, not on the evidence alone.

\begin{table}[t]
\centering\small
\begin{tabular}{lllcccc}
\toprule
Model & Budget & Task & TF-IDF & DeBERTa (q+added) & DeBERTa (+context) & Probe \\
\midrule
Qwen2.5-7B & reduced budget & sufficiency & 0.623 & 0.723 & 0.900 & 0.948 \\
Qwen2.5-7B & reduced budget & sufficiency matched & 0.623 & 0.758 & 0.823 & 0.911 \\
Qwen2.5-7B & reduced budget & uptake & 0.574 & 0.569 & 0.557 & 0.775 \\
Qwen2.5-7B & full budget & sufficiency & 0.611 & 0.723 & 0.972 & 0.948 \\
Qwen2.5-7B & full budget & uptake & 0.574 & 0.625 & 0.630 & 0.775 \\
Llama-3.1-8B & reduced budget & sufficiency & 0.625 & 0.715 & 0.888 & 0.952 \\
Llama-3.1-8B & reduced budget & uptake & 0.463 & 0.466 & 0.517 & 0.783 \\
Mistral-7B & reduced budget & sufficiency & 0.625 & 0.730 & 0.920 & 0.949 \\
Mistral-7B & reduced budget & uptake & 0.536 & 0.573 & 0.604 & 0.771 \\
\bottomrule
\end{tabular}
\caption{Text-only classifiers on the same increments and splits versus the activation probe (document split).}
\label{tab:text}
\end{table}


